% ---
% RESUMOS
% ---

% resumo em português
\setlength{\absparsep}{18pt} % ajusta o espaçamento dos parágrafos do resumo
\begin{resumo}
Este trabalho desenvolve um sistema integrado de aquisição para a bancada de ensaios de conjuntos motopropulsores com motores BLDC do projeto Albatroz Aerodesign. Uma placa de medição, baseada no microcontrolador CH32V307, adquire as tensões e correntes de fase do motor, aciona o controlador do motor e os relés de potência e concentra a segurança, comunicando-se por Ethernet e UDP com um \textit{software} de supervisão. Para estimar a velocidade angular e o torque no eixo sem sensores mecânicos, implementa-se, embarcado na placa, um observador por modos deslizantes (SMO) associado a um PLL e à reconstrução analítica da força contraeletromotriz. O sistema permite comparar as grandezas medidas com as estimadas pelo observador, servindo como ferramenta de caracterização de conjuntos motopropulsores.

 \textbf{Palavras-chave}: Bancada de Ensaios; Motor BLDC; Observador por Modos Deslizantes; Aquisição de Dados; Sistema \textit{Sensorless}; Propulsão Elétrica.
\end{resumo}
